% DRAFT: paper §6.3 — combined E5 null finding (Arm 2 + Arm 3).
% Updated 2026-06-14 after Arm 3 (qwen-7B r=64) also fired NO-GO and
% the mechanical-forcing math made the unreachability of the
% floor-positive regime explicit at practical LoRA ranks.

\subsection{Why the floor formula is operationally vacuous on real LoRA fine-tunes}
\label{subsec:exp-e5-null}

The lower bound of \S\ref{sec:lower_bound} predicts a non-zero
worst-task distortion floor of $B^2(1 - d_{\mathrm{eff}}/(Tr))$
whenever the task subspaces partially overlap
($d_{\mathrm{eff}} < Tr$). On the sixteen original adapters of
\S\ref{subsec:exp-lora} we measured $d_{\mathrm{eff}} = Tr$ in every
attention projection of every base model, so the predicted floor is
identically zero. A natural reviewer concern: is this an unlucky
choice of task corpora, or fine-tuning hyperparameters, that
happens to land on the saturation boundary? We tested two
operational routes to break the saturation and report the results
together.

\paragraph{Arm 2: data-mixture overlap.}
On Qwen-2.5-7B we trained $4$ task adapters under three
data-mixture levels: at mixing parameter $\alpha \in \{0, 0.5, 0.9\}$
the training data for the target task is composed of a fraction
$\alpha$ from a fixed shared pool (uniformly sampled across all
four task corpora) and a fraction $(1 - \alpha)$ from the target
task's own training set. At $\alpha = 0.9$ four adapters share
$90\%$ of their training examples. Per-layer hard
$d_{\mathrm{eff}}/(Tr)$ stayed saturated at $1.000$ across all $112$
attention projections at every $\alpha$ (Table~\ref{tab:e5-null}).
Soft (participation-ratio) $d_{\mathrm{eff}}/(Tr)$ drifted by
$2\%$ ($0.257 \to 0.251$), an amount well inside the seed-noise
scale of \S\ref{subsec:exp-multiseed}. \emph{Real fine-tuning
resists subspace overlap}: the $V_t = \text{top-}r$ right-singular
basis of $\Delta_t$ is set by the per-task answer-distribution
loss, not by the inputs' raw geometry, so shared training inputs
do not induce shared LoRA subspaces.

\paragraph{Arm 3: geometric forcing at $r=64$.}
On the same Qwen-2.5-7B base we then trained the four task
adapters at rank $r = 64$ (vs.\ $r = 16$ in the main matrix),
giving $T \cdot r = 256$. Hard $d_{\mathrm{eff}}/(Tr)$: mean $0.997$,
min $0.867$, $0/112$ layers below the pre-specified $0.8$
threshold. The single most-bottlenecked layer drops to
$\approx 0.867$, but the floor formula still evaluates to zero in
the majority of layers and is operationally negligible across the
network. (A cross-architecture run on
Llama-3.2-3B at $r=64$ would have been an empirical confirmation,
but the mechanical bound below makes the question purely
mathematical: with $T = 4$ and $\mathrm{in\_dim} = 3072$, Llama-3.2-3B
would need $r > 768$ for $d_{\mathrm{eff}}/(Tr) < 1$ to be even
mathematically attainable, so any practical-rank result must be at
the saturation point.)

\paragraph{The mechanical-forcing barrier.}
The reason both arms fail to break saturation is now precise.
Generically, $d_{\mathrm{eff}} = \mathrm{rank}(\sum_t P_{V_t}) =
\min(\mathrm{in\_dim},\, T \cdot r)$. So
$d_{\mathrm{eff}}/(Tr) < 1$ requires $T \cdot r > \mathrm{in\_dim}$,
i.e.\ $r > \mathrm{in\_dim}/T$. Per-model: this is
$r > 1024$ on Llama-3.1-8B, $r > 896$ on Qwen-2.5-7B,
$r > 768$ on Llama-3.2-3B, and similar for Mistral and Yi. These
are all well outside the practical LoRA regime (which uses
$r \leq 64$ in nearly all production deployments). At
$T = 4$ tasks, mechanically forcing the floor-positive regime
would require LoRAs of rank comparable to the model's hidden
dimension, defeating the parameter-efficiency that motivates LoRA
in the first place.

\begin{table}[t]
\centering
\caption{Subspace metrics at Qwen-2.5-7B's $112$ attention
projections under the two E5 arms. Arm 2: data-mixture overlap at
$\alpha = 0.9$. Arm 3: geometric forcing at $r=64$. Neither breaks
the $d_{\mathrm{eff}}/(Tr) = 1$ saturation in any majority sense.
The hard variant is the literal bound; the soft (participation-ratio)
variant is included as a finer measurement of angular alignment.}
\label{tab:e5-null}
\begin{tabular}{lrrrrr}
\hline
Condition & hard mean & hard min & soft mean & layers $< 0.8$ \\
\hline
Arm 2, $\alpha = 0.0$  & $1.000$ & $1.000$ & $0.257$ & $0/112$ \\
Arm 2, $\alpha = 0.5$  & $1.000$ & $1.000$ & $0.253$ & $0/112$ \\
Arm 2, $\alpha = 0.9$  & $1.000$ & $1.000$ & $0.251$ & $0/112$ \\
Arm 3, $r = 64$        & $0.997$ & $0.867$ & $0.253$ & $0/112$ \\
\hline
\end{tabular}
\end{table}

\paragraph{Interpretation: the floor formula is real but operationally vacuous.}
The lower bound continues to apply --- it is an information-theoretic
floor that always holds --- but its non-vacuous regime sits outside
the practical LoRA setting in any natural way we could probe. Real
fine-tuning consistently lands at $d_{\mathrm{eff}} = Tr$ for any
reasonable $r$, and the term $1 - d_{\mathrm{eff}}/(Tr)$ in the
floor expression therefore vanishes operationally. This is a
\emph{positive} finding for the achievability program: the
\emph{algorithmic} regime, where worst-task excess is dominated by
construction-side suboptimality rather than the irreducible floor,
is the regime that all real LoRA fine-tunes live in. The
regularized achievability salvage of
\S\ref{subsec:exp-rd-encoder} ---
$(\bar H + \lambda I)^{-1}$ with $\lambda \in [0.05, 0.13]$ ---
sits squarely on the universally applicable side of the bound, and
gives the practical method that practitioners can use.

We do not rule out a deliberately engineered floor-positive regime
(e.g., shared adapter modules across tasks, hard-coded subspace
constraints, or extreme high-rank LoRA approaching dense
fine-tuning). The construction of such regimes is left as future
work; the present empirical finding is that they do not arise
naturally from data, hyperparameter, or rank choices a
practitioner would make.

%%% END DRAFT %%%
